% Übernahme und Rückgabe der Arme (Überblendung per Gewicht), schematisch. Keine Messwerte, keine Skalen.
% Eigene Darstellung (Art E, Kap. 6.3.3). Keine Quelle in der Caption.
% Basis: G1@6a40e47 g1pilot/docs/20-T-Arm-Manipulation.md, Abschnitte "rt/arm_sdk-Gewichtsrampe" (Gewicht 0..1,
%   Rampe über arm_weight_ramp_up_s, real 2,0 s, "beim DISABLE symmetrisch zurück") und "Übergabe an den
%   Onboard-Regler" (_record_onboard_pose merkt die gemessene Armpose; DISABLE -> geplante Bahn in die gemerkte
%   Pose, erst bei Ankunft Rampe 1->0; Zeitgrenze release_timeout_s 12 s, hier nicht gezeichnet);
%   docs/50-A-Echter-Roboter.md E ("Übernehmen ohne Bewegung": gemessene Pose = Pose des Onboard-Reglers,
%   kommandiert genau diese, Gewicht 0->1 über 2,0 s).
% a) Übernehmen: Gewicht 0 -> 1, Abstand der kommandierten Pose zur gemerkten Onboard-Pose bleibt 0.
% b) Zurückgeben: Gewicht bleibt 1 während der geplanten Rückfahrt (Abstand sinkt auf 0), danach Rampe 1 -> 0.
% Nicht gezeichnet: Not-Halt (Gewicht bleibt 1, Steifigkeit 0) und Programmende (kurze Rückfahrt, dann Rampe).
% Stil: ba-grafik.tex. Einbinden: % Übernahme und Rückgabe der Arme (Überblendung per Gewicht), schematisch. Keine Messwerte, keine Skalen.
% Eigene Darstellung (Art E, Kap. 6.3.3). Keine Quelle in der Caption.
% Basis: G1@6a40e47 g1pilot/docs/20-T-Arm-Manipulation.md, Abschnitte "rt/arm_sdk-Gewichtsrampe" (Gewicht 0..1,
%   Rampe über arm_weight_ramp_up_s, real 2,0 s, "beim DISABLE symmetrisch zurück") und "Übergabe an den
%   Onboard-Regler" (_record_onboard_pose merkt die gemessene Armpose; DISABLE -> geplante Bahn in die gemerkte
%   Pose, erst bei Ankunft Rampe 1->0; Zeitgrenze release_timeout_s 12 s, hier nicht gezeichnet);
%   docs/50-A-Echter-Roboter.md E ("Übernehmen ohne Bewegung": gemessene Pose = Pose des Onboard-Reglers,
%   kommandiert genau diese, Gewicht 0->1 über 2,0 s).
% a) Übernehmen: Gewicht 0 -> 1, Abstand der kommandierten Pose zur gemerkten Onboard-Pose bleibt 0.
% b) Zurückgeben: Gewicht bleibt 1 während der geplanten Rückfahrt (Abstand sinkt auf 0), danach Rampe 1 -> 0.
% Nicht gezeichnet: Not-Halt (Gewicht bleibt 1, Steifigkeit 0) und Programmende (kurze Rückfahrt, dann Rampe).
% Stil: ba-grafik.tex. Einbinden: % Übernahme und Rückgabe der Arme (Überblendung per Gewicht), schematisch. Keine Messwerte, keine Skalen.
% Eigene Darstellung (Art E, Kap. 6.3.3). Keine Quelle in der Caption.
% Basis: G1@6a40e47 g1pilot/docs/20-T-Arm-Manipulation.md, Abschnitte "rt/arm_sdk-Gewichtsrampe" (Gewicht 0..1,
%   Rampe über arm_weight_ramp_up_s, real 2,0 s, "beim DISABLE symmetrisch zurück") und "Übergabe an den
%   Onboard-Regler" (_record_onboard_pose merkt die gemessene Armpose; DISABLE -> geplante Bahn in die gemerkte
%   Pose, erst bei Ankunft Rampe 1->0; Zeitgrenze release_timeout_s 12 s, hier nicht gezeichnet);
%   docs/50-A-Echter-Roboter.md E ("Übernehmen ohne Bewegung": gemessene Pose = Pose des Onboard-Reglers,
%   kommandiert genau diese, Gewicht 0->1 über 2,0 s).
% a) Übernehmen: Gewicht 0 -> 1, Abstand der kommandierten Pose zur gemerkten Onboard-Pose bleibt 0.
% b) Zurückgeben: Gewicht bleibt 1 während der geplanten Rückfahrt (Abstand sinkt auf 0), danach Rampe 1 -> 0.
% Nicht gezeichnet: Not-Halt (Gewicht bleibt 1, Steifigkeit 0) und Programmende (kurze Rückfahrt, dann Rampe).
% Stil: ba-grafik.tex. Einbinden: % Übernahme und Rückgabe der Arme (Überblendung per Gewicht), schematisch. Keine Messwerte, keine Skalen.
% Eigene Darstellung (Art E, Kap. 6.3.3). Keine Quelle in der Caption.
% Basis: G1@6a40e47 g1pilot/docs/20-T-Arm-Manipulation.md, Abschnitte "rt/arm_sdk-Gewichtsrampe" (Gewicht 0..1,
%   Rampe über arm_weight_ramp_up_s, real 2,0 s, "beim DISABLE symmetrisch zurück") und "Übergabe an den
%   Onboard-Regler" (_record_onboard_pose merkt die gemessene Armpose; DISABLE -> geplante Bahn in die gemerkte
%   Pose, erst bei Ankunft Rampe 1->0; Zeitgrenze release_timeout_s 12 s, hier nicht gezeichnet);
%   docs/50-A-Echter-Roboter.md E ("Übernehmen ohne Bewegung": gemessene Pose = Pose des Onboard-Reglers,
%   kommandiert genau diese, Gewicht 0->1 über 2,0 s).
% a) Übernehmen: Gewicht 0 -> 1, Abstand der kommandierten Pose zur gemerkten Onboard-Pose bleibt 0.
% b) Zurückgeben: Gewicht bleibt 1 während der geplanten Rückfahrt (Abstand sinkt auf 0), danach Rampe 1 -> 0.
% Nicht gezeichnet: Not-Halt (Gewicht bleibt 1, Steifigkeit 0) und Programmende (kurze Rückfahrt, dann Rampe).
% Stil: ba-grafik.tex. Einbinden: \input{Bilder/tikz/armuebergabe} in einer figure-Umgebung. Nicht skalieren.
\begin{tikzpicture}[ba]
  \tikzset{
    achse/.style={ba normal, ->, color=baGrau},
    gew/.style={color=baGruen, line width=1.6pt},
    abst/.style={color=baBlau, line width=1.6pt, dash pattern=on 5pt off 2pt},
    ereig/.style={ba hilfslinie},
    elabel/.style={font=\footnotesize, anchor=south, align=center},
    klein/.style={font=\footnotesize, text=baGrau},
    tick/.style={font=\footnotesize, anchor=east},
  }
  % Layout für 12pt-Grundschrift der Arbeit geprüft (Achsen 5,4 cm, Teilbilder weit genug auseinander, damit
  % die Achsenbeschriftung von b) nicht in a) ragt).
  % ======================== a) Übernehmen ========================
  \begin{scope}[shift={(1.35,0)}]
    % Gewicht
    \draw[achse] (0,4.1) -- (0,6.0);
    \draw[achse] (0,4.1) -- (5.4,4.1);
    \node[tick] at (-0.1,4.1) {0};
    \node[tick] at (-0.1,5.7) {1};
    \draw[ba duenn, color=baGrau] (-0.05,5.7) -- (0.05,5.7);
    \node[font=\footnotesize, rotate=90, anchor=south] at (-0.75,5.05) {Gewicht};
    \draw[gew] (0,4.1) -- (1.0,4.1) -- (3.0,5.7) -- (5.2,5.7);
    \node[klein, anchor=west] at (2.6,4.5) {Rampe über 2\,s};
    % Abstand
    \draw[achse] (0,1.0) -- (0,2.9);
    \draw[achse] (0,1.0) -- (5.4,1.0) node[font=\footnotesize, anchor=north east] {Zeit};
    \node[tick] at (-0.1,1.0) {0};
    \node[font=\footnotesize, rotate=90, anchor=south, align=center] at (-0.75,1.95) {Abstand zur\\Onboard-Pose};
    \draw[abst] (0,1.0) -- (5.2,1.0);
    \node[klein, anchor=west, align=left] at (1.3,1.65) {kommandiert wird die\\gemessene Pose};
    % Ereignis
    \draw[ereig] (1.0,0.85) -- (1.0,6.2);
    \node[elabel] at (1.0,6.25) {Übernehmen};
  \end{scope}
  % Trenner
  \draw[ba trenner] (7.1,0.0) -- (7.1,6.9); % mittig zwischen Achsende a) und Beschriftung b)
  % ======================== b) Zurückgeben ========================
  \begin{scope}[shift={(9.0,0)}]
    % Gewicht
    \draw[achse] (0,4.1) -- (0,6.0);
    \draw[achse] (0,4.1) -- (5.4,4.1);
    \node[tick] at (-0.1,4.1) {0};
    \node[tick] at (-0.1,5.7) {1};
    \draw[ba duenn, color=baGrau] (-0.05,5.7) -- (0.05,5.7);
    \node[font=\footnotesize, rotate=90, anchor=south] at (-0.75,5.05) {Gewicht};
    \draw[gew] (0,5.7) -- (3.3,5.7) -- (5.0,4.1);
    \node[klein, anchor=west] at (4.3,5.3) {Rampe};
    % Abstand
    \draw[achse] (0,1.0) -- (0,2.9);
    \draw[achse] (0,1.0) -- (5.4,1.0) node[font=\footnotesize, anchor=north east] {Zeit};
    \node[tick] at (-0.1,1.0) {0};
    \node[font=\footnotesize, rotate=90, anchor=south, align=center] at (-0.75,1.95) {Abstand zur\\Onboard-Pose};
    \draw[abst] (0,2.6) -- (0.8,2.6) .. controls (1.7,2.6) and (2.2,1.0) .. (3.3,1.0) -- (5.0,1.0);
    \node[klein, anchor=west, align=left] at (3.4,2.05) {geplante,\\langsame\\Rückfahrt};
    % Ereignisse
    \draw[ereig] (0.8,0.85) -- (0.8,6.2);
    \node[elabel] at (0.8,6.25) {Zurückgeben};
    \draw[ereig] (3.3,0.85) -- (3.3,6.2);
    \node[elabel] at (3.3,6.25) {Pose erreicht};
  \end{scope}
  % Teiltitel
  \node[ba teiltitel] at (4.05,0.2) {a) Übernehmen};
  \node[ba teiltitel] at (11.7,0.2) {b) Zurückgeben};
  % Legende
  \draw[gew] (1.35,-0.75) -- (2.15,-0.75);
  \node[klein, anchor=west] at (2.25,-0.75) {Gewicht der eigenen Armregelung};
  \draw[abst] (9.0,-0.75) -- (9.8,-0.75);
  \node[klein, anchor=west] at (9.9,-0.75) {Abstand zur Onboard-Pose};
\end{tikzpicture}
 in einer figure-Umgebung. Nicht skalieren.
\begin{tikzpicture}[ba]
  \tikzset{
    achse/.style={ba normal, ->, color=baGrau},
    gew/.style={color=baGruen, line width=1.6pt},
    abst/.style={color=baBlau, line width=1.6pt, dash pattern=on 5pt off 2pt},
    ereig/.style={ba hilfslinie},
    elabel/.style={font=\footnotesize, anchor=south, align=center},
    klein/.style={font=\footnotesize, text=baGrau},
    tick/.style={font=\footnotesize, anchor=east},
  }
  % Layout für 12pt-Grundschrift der Arbeit geprüft (Achsen 5,4 cm, Teilbilder weit genug auseinander, damit
  % die Achsenbeschriftung von b) nicht in a) ragt).
  % ======================== a) Übernehmen ========================
  \begin{scope}[shift={(1.35,0)}]
    % Gewicht
    \draw[achse] (0,4.1) -- (0,6.0);
    \draw[achse] (0,4.1) -- (5.4,4.1);
    \node[tick] at (-0.1,4.1) {0};
    \node[tick] at (-0.1,5.7) {1};
    \draw[ba duenn, color=baGrau] (-0.05,5.7) -- (0.05,5.7);
    \node[font=\footnotesize, rotate=90, anchor=south] at (-0.75,5.05) {Gewicht};
    \draw[gew] (0,4.1) -- (1.0,4.1) -- (3.0,5.7) -- (5.2,5.7);
    \node[klein, anchor=west] at (2.6,4.5) {Rampe über 2\,s};
    % Abstand
    \draw[achse] (0,1.0) -- (0,2.9);
    \draw[achse] (0,1.0) -- (5.4,1.0) node[font=\footnotesize, anchor=north east] {Zeit};
    \node[tick] at (-0.1,1.0) {0};
    \node[font=\footnotesize, rotate=90, anchor=south, align=center] at (-0.75,1.95) {Abstand zur\\Onboard-Pose};
    \draw[abst] (0,1.0) -- (5.2,1.0);
    \node[klein, anchor=west, align=left] at (1.3,1.65) {kommandiert wird die\\gemessene Pose};
    % Ereignis
    \draw[ereig] (1.0,0.85) -- (1.0,6.2);
    \node[elabel] at (1.0,6.25) {Übernehmen};
  \end{scope}
  % Trenner
  \draw[ba trenner] (7.1,0.0) -- (7.1,6.9); % mittig zwischen Achsende a) und Beschriftung b)
  % ======================== b) Zurückgeben ========================
  \begin{scope}[shift={(9.0,0)}]
    % Gewicht
    \draw[achse] (0,4.1) -- (0,6.0);
    \draw[achse] (0,4.1) -- (5.4,4.1);
    \node[tick] at (-0.1,4.1) {0};
    \node[tick] at (-0.1,5.7) {1};
    \draw[ba duenn, color=baGrau] (-0.05,5.7) -- (0.05,5.7);
    \node[font=\footnotesize, rotate=90, anchor=south] at (-0.75,5.05) {Gewicht};
    \draw[gew] (0,5.7) -- (3.3,5.7) -- (5.0,4.1);
    \node[klein, anchor=west] at (4.3,5.3) {Rampe};
    % Abstand
    \draw[achse] (0,1.0) -- (0,2.9);
    \draw[achse] (0,1.0) -- (5.4,1.0) node[font=\footnotesize, anchor=north east] {Zeit};
    \node[tick] at (-0.1,1.0) {0};
    \node[font=\footnotesize, rotate=90, anchor=south, align=center] at (-0.75,1.95) {Abstand zur\\Onboard-Pose};
    \draw[abst] (0,2.6) -- (0.8,2.6) .. controls (1.7,2.6) and (2.2,1.0) .. (3.3,1.0) -- (5.0,1.0);
    \node[klein, anchor=west, align=left] at (3.4,2.05) {geplante,\\langsame\\Rückfahrt};
    % Ereignisse
    \draw[ereig] (0.8,0.85) -- (0.8,6.2);
    \node[elabel] at (0.8,6.25) {Zurückgeben};
    \draw[ereig] (3.3,0.85) -- (3.3,6.2);
    \node[elabel] at (3.3,6.25) {Pose erreicht};
  \end{scope}
  % Teiltitel
  \node[ba teiltitel] at (4.05,0.2) {a) Übernehmen};
  \node[ba teiltitel] at (11.7,0.2) {b) Zurückgeben};
  % Legende
  \draw[gew] (1.35,-0.75) -- (2.15,-0.75);
  \node[klein, anchor=west] at (2.25,-0.75) {Gewicht der eigenen Armregelung};
  \draw[abst] (9.0,-0.75) -- (9.8,-0.75);
  \node[klein, anchor=west] at (9.9,-0.75) {Abstand zur Onboard-Pose};
\end{tikzpicture}
 in einer figure-Umgebung. Nicht skalieren.
\begin{tikzpicture}[ba]
  \tikzset{
    achse/.style={ba normal, ->, color=baGrau},
    gew/.style={color=baGruen, line width=1.6pt},
    abst/.style={color=baBlau, line width=1.6pt, dash pattern=on 5pt off 2pt},
    ereig/.style={ba hilfslinie},
    elabel/.style={font=\footnotesize, anchor=south, align=center},
    klein/.style={font=\footnotesize, text=baGrau},
    tick/.style={font=\footnotesize, anchor=east},
  }
  % Layout für 12pt-Grundschrift der Arbeit geprüft (Achsen 5,4 cm, Teilbilder weit genug auseinander, damit
  % die Achsenbeschriftung von b) nicht in a) ragt).
  % ======================== a) Übernehmen ========================
  \begin{scope}[shift={(1.35,0)}]
    % Gewicht
    \draw[achse] (0,4.1) -- (0,6.0);
    \draw[achse] (0,4.1) -- (5.4,4.1);
    \node[tick] at (-0.1,4.1) {0};
    \node[tick] at (-0.1,5.7) {1};
    \draw[ba duenn, color=baGrau] (-0.05,5.7) -- (0.05,5.7);
    \node[font=\footnotesize, rotate=90, anchor=south] at (-0.75,5.05) {Gewicht};
    \draw[gew] (0,4.1) -- (1.0,4.1) -- (3.0,5.7) -- (5.2,5.7);
    \node[klein, anchor=west] at (2.6,4.5) {Rampe über 2\,s};
    % Abstand
    \draw[achse] (0,1.0) -- (0,2.9);
    \draw[achse] (0,1.0) -- (5.4,1.0) node[font=\footnotesize, anchor=north east] {Zeit};
    \node[tick] at (-0.1,1.0) {0};
    \node[font=\footnotesize, rotate=90, anchor=south, align=center] at (-0.75,1.95) {Abstand zur\\Onboard-Pose};
    \draw[abst] (0,1.0) -- (5.2,1.0);
    \node[klein, anchor=west, align=left] at (1.3,1.65) {kommandiert wird die\\gemessene Pose};
    % Ereignis
    \draw[ereig] (1.0,0.85) -- (1.0,6.2);
    \node[elabel] at (1.0,6.25) {Übernehmen};
  \end{scope}
  % Trenner
  \draw[ba trenner] (7.1,0.0) -- (7.1,6.9); % mittig zwischen Achsende a) und Beschriftung b)
  % ======================== b) Zurückgeben ========================
  \begin{scope}[shift={(9.0,0)}]
    % Gewicht
    \draw[achse] (0,4.1) -- (0,6.0);
    \draw[achse] (0,4.1) -- (5.4,4.1);
    \node[tick] at (-0.1,4.1) {0};
    \node[tick] at (-0.1,5.7) {1};
    \draw[ba duenn, color=baGrau] (-0.05,5.7) -- (0.05,5.7);
    \node[font=\footnotesize, rotate=90, anchor=south] at (-0.75,5.05) {Gewicht};
    \draw[gew] (0,5.7) -- (3.3,5.7) -- (5.0,4.1);
    \node[klein, anchor=west] at (4.3,5.3) {Rampe};
    % Abstand
    \draw[achse] (0,1.0) -- (0,2.9);
    \draw[achse] (0,1.0) -- (5.4,1.0) node[font=\footnotesize, anchor=north east] {Zeit};
    \node[tick] at (-0.1,1.0) {0};
    \node[font=\footnotesize, rotate=90, anchor=south, align=center] at (-0.75,1.95) {Abstand zur\\Onboard-Pose};
    \draw[abst] (0,2.6) -- (0.8,2.6) .. controls (1.7,2.6) and (2.2,1.0) .. (3.3,1.0) -- (5.0,1.0);
    \node[klein, anchor=west, align=left] at (3.4,2.05) {geplante,\\langsame\\Rückfahrt};
    % Ereignisse
    \draw[ereig] (0.8,0.85) -- (0.8,6.2);
    \node[elabel] at (0.8,6.25) {Zurückgeben};
    \draw[ereig] (3.3,0.85) -- (3.3,6.2);
    \node[elabel] at (3.3,6.25) {Pose erreicht};
  \end{scope}
  % Teiltitel
  \node[ba teiltitel] at (4.05,0.2) {a) Übernehmen};
  \node[ba teiltitel] at (11.7,0.2) {b) Zurückgeben};
  % Legende
  \draw[gew] (1.35,-0.75) -- (2.15,-0.75);
  \node[klein, anchor=west] at (2.25,-0.75) {Gewicht der eigenen Armregelung};
  \draw[abst] (9.0,-0.75) -- (9.8,-0.75);
  \node[klein, anchor=west] at (9.9,-0.75) {Abstand zur Onboard-Pose};
\end{tikzpicture}
 in einer figure-Umgebung. Nicht skalieren.
\begin{tikzpicture}[ba]
  \tikzset{
    achse/.style={ba normal, ->, color=baGrau},
    gew/.style={color=baGruen, line width=1.6pt},
    abst/.style={color=baBlau, line width=1.6pt, dash pattern=on 5pt off 2pt},
    ereig/.style={ba hilfslinie},
    elabel/.style={font=\footnotesize, anchor=south, align=center},
    klein/.style={font=\footnotesize, text=baGrau},
    tick/.style={font=\footnotesize, anchor=east},
  }
  % Layout für 12pt-Grundschrift der Arbeit geprüft (Achsen 5,4 cm, Teilbilder weit genug auseinander, damit
  % die Achsenbeschriftung von b) nicht in a) ragt).
  % ======================== a) Übernehmen ========================
  \begin{scope}[shift={(1.35,0)}]
    % Gewicht
    \draw[achse] (0,4.1) -- (0,6.0);
    \draw[achse] (0,4.1) -- (5.4,4.1);
    \node[tick] at (-0.1,4.1) {0};
    \node[tick] at (-0.1,5.7) {1};
    \draw[ba duenn, color=baGrau] (-0.05,5.7) -- (0.05,5.7);
    \node[font=\footnotesize, rotate=90, anchor=south] at (-0.75,5.05) {Gewicht};
    \draw[gew] (0,4.1) -- (1.0,4.1) -- (3.0,5.7) -- (5.2,5.7);
    \node[klein, anchor=west] at (2.6,4.5) {Rampe über 2\,s};
    % Abstand
    \draw[achse] (0,1.0) -- (0,2.9);
    \draw[achse] (0,1.0) -- (5.4,1.0) node[font=\footnotesize, anchor=north east] {Zeit};
    \node[tick] at (-0.1,1.0) {0};
    \node[font=\footnotesize, rotate=90, anchor=south, align=center] at (-0.75,1.95) {Abstand zur\\Onboard-Pose};
    \draw[abst] (0,1.0) -- (5.2,1.0);
    \node[klein, anchor=west, align=left] at (1.3,1.65) {kommandiert wird die\\gemessene Pose};
    % Ereignis
    \draw[ereig] (1.0,0.85) -- (1.0,6.2);
    \node[elabel] at (1.0,6.25) {Übernehmen};
  \end{scope}
  % Trenner
  \draw[ba trenner] (7.1,0.0) -- (7.1,6.9); % mittig zwischen Achsende a) und Beschriftung b)
  % ======================== b) Zurückgeben ========================
  \begin{scope}[shift={(9.0,0)}]
    % Gewicht
    \draw[achse] (0,4.1) -- (0,6.0);
    \draw[achse] (0,4.1) -- (5.4,4.1);
    \node[tick] at (-0.1,4.1) {0};
    \node[tick] at (-0.1,5.7) {1};
    \draw[ba duenn, color=baGrau] (-0.05,5.7) -- (0.05,5.7);
    \node[font=\footnotesize, rotate=90, anchor=south] at (-0.75,5.05) {Gewicht};
    \draw[gew] (0,5.7) -- (3.3,5.7) -- (5.0,4.1);
    \node[klein, anchor=west] at (4.3,5.3) {Rampe};
    % Abstand
    \draw[achse] (0,1.0) -- (0,2.9);
    \draw[achse] (0,1.0) -- (5.4,1.0) node[font=\footnotesize, anchor=north east] {Zeit};
    \node[tick] at (-0.1,1.0) {0};
    \node[font=\footnotesize, rotate=90, anchor=south, align=center] at (-0.75,1.95) {Abstand zur\\Onboard-Pose};
    \draw[abst] (0,2.6) -- (0.8,2.6) .. controls (1.7,2.6) and (2.2,1.0) .. (3.3,1.0) -- (5.0,1.0);
    \node[klein, anchor=west, align=left] at (3.4,2.05) {geplante,\\langsame\\Rückfahrt};
    % Ereignisse
    \draw[ereig] (0.8,0.85) -- (0.8,6.2);
    \node[elabel] at (0.8,6.25) {Zurückgeben};
    \draw[ereig] (3.3,0.85) -- (3.3,6.2);
    \node[elabel] at (3.3,6.25) {Pose erreicht};
  \end{scope}
  % Teiltitel
  \node[ba teiltitel] at (4.05,0.2) {a) Übernehmen};
  \node[ba teiltitel] at (11.7,0.2) {b) Zurückgeben};
  % Legende
  \draw[gew] (1.35,-0.75) -- (2.15,-0.75);
  \node[klein, anchor=west] at (2.25,-0.75) {Gewicht der eigenen Armregelung};
  \draw[abst] (9.0,-0.75) -- (9.8,-0.75);
  \node[klein, anchor=west] at (9.9,-0.75) {Abstand zur Onboard-Pose};
\end{tikzpicture}
